% !PRINTED 265
\weberpage{265}
%<<CONT
negativem Zeichen genommen, von $r$ nach dem Modul $\frac{1}{2}(p-1)$
congruent sind. Man erhält daher alle von einander und von
$\pm 2$ verschiedenen Werthe dieses Ausdruckes, wenn man
\weq{10}{r = 1, 2, \cdots \frac{p-3}{2}}
setzt. Da hier für jeden Exponenten $k$
\[
S^k = \begin{pmatrix} \sigma^k, & 0 \\ 0, & \sigma^{-k} \end{pmatrix}
\]
ist, so ist der Grad von $S$, und also auch von jeder Substitution
der II$^{\mathrm{ten}}$ Art entweder $\frac{1}{2}(p-1)$ oder ein Theiler davon, je
nachdem $r$ relativ prim zu $\frac{1}{2}(p-1)$ ist oder nicht.

Um die Anzahl dieser Substitutionen zu bestimmen, verfahren
 wir wie oben. Ist zunächst $c = 0$, so ergiebt sich
\[
a = \gamma^{\pm r(p+1)}, \quad d = \gamma^{\mp r(p+1)},
\]
und $b$ kann $p$ Werthe haben. Die Zahl $r$ hat die Werthe (10),
und also ergeben sich hiernach $p\,(p-3)$ Substitutionen. Ist
dann $c$ von Null verschieden, so ist $p\,(p-1)$ die Anzahl der
verschiedenen Annahmen über $a, c$. Ist $a$ angenommen, so können
wir $d$ aus (9) bestimmen und
\[
b = -c^{-1} + a\,d\,c^{-1},
\]
setzen, woraus wir $\frac{1}{2}\,p\,(p-1)(p-3)$ Bestimmungen erhalten,
also mit den für $c = 0$ gezählten zusammen $\frac{1}{2}\,p\,(p-3)(p+1)$.
Auch diese Zahl ist noch zu halbiren, so dass wir
\[
\frac{1}{4}\,p\,(p-3)(p+1)
\]
Substitutionen der II$^{\mathrm{ten}}$ Art erhalten.

III.~$(a+d)^2 - 4$ quadratischer Nichtrest von $p$. In diesem
Falle ist $\sigma$ nicht reell, aber mit seinem reciproken Werth
conjugirt. Also ist nach \S.\,67, (4) $\sigma^p = \sigma^{-1}$, und folglich
\[
\sigma^{\frac{p+1}{2}} = \pm 1.
\]
Daraus ergiebt sich, dass der Grad einer solchen Substitution
 $\frac{1}{2}(p+1)$ oder ein Theiler dieser Zahl ist.
In diesem dritten Falle setzen wir
\weq{11}{\begin{aligned}
&\sigma = \gamma^{r(p-1)}, \quad a+d = \gamma^{r(p-1)} + \gamma^{-r(p-1)}\\
&r = 1, 2, \ldots, \tfrac{1}{2}(p-1).
\end{aligned}}
Hier kann der Fall $c = 0$ oder $b = 0$ nicht vorkommen,
weil sich aus $a\,d = 1$ ergeben würde, dass $a = \sigma^{\pm 1}$, $d = \sigma^{\mp 1}$
%CONT>>
